% =====================================================================
% AWPL letter
% Closed-form peak spatial-average SAR over the IEEE/IEC 10 g cube
% above 6 GHz, derived from the surface absorbed power density.
% Author: Robin Wydaeghe
% Target: IEEE Antennas and Wireless Propagation Letters (AWPL)
% =====================================================================

\documentclass[journal,twocolumn,10pt]{IEEEtran}

\usepackage[utf8]{inputenc}
\usepackage[T1]{fontenc}
\usepackage{lmodern}
\usepackage{microtype}

\usepackage{amsmath,amssymb,amsthm,mathtools}
\usepackage{bm}

\usepackage{graphicx}
\graphicspath{{figures/psSAR/}{figures/}{../../theory/figures/}{../../validation/}}
\usepackage{booktabs}
\usepackage{array}
\usepackage{caption}

\usepackage{cite}
\usepackage{xcolor}
\usepackage[colorlinks=true,linkcolor=blue,citecolor=blue,urlcolor=blue]{hyperref}
\usepackage[capitalize]{cleveref}

\theoremstyle{plain}
\newtheorem{proposition}{Proposition}
\newtheorem{corollary}[proposition]{Corollary}
\theoremstyle{remark}
\newtheorem{remark}[proposition]{Remark}

\crefname{proposition}{Proposition}{Propositions}
\Crefname{proposition}{Proposition}{Propositions}
\crefname{remark}{Remark}{Remarks}
\Crefname{remark}{Remark}{Remarks}

% --- math macros (kept consistent with the sister papers) ---
\newcommand{\khat}{\hat{\bm{k}}}
\newcommand{\nhat}{\hat{\bm{n}}}
\newcommand{\rr}{\mathbf{r}}
\newcommand{\Sinc}{S_{\mathrm{inc}}}
\newcommand{\Sab}{S_{\mathrm{ab}}}
\newcommand{\Aproj}{A_\perp}
\newcommand{\rhom}{\rho_{\mathrm{m}}}
\newcommand{\ntilde}{\tilde{n}}
\newcommand{\diff}{\mathrm{d}}
\newcommand{\pospart}[1]{\left(#1\right)_{+}}
\DeclareMathOperator{\RE}{Re}

\begin{document}

\title{A closed-form peak spatial-average {SAR} over the
{IEEE/IEC} 10\,g cube above 6\,{GHz}}

\author{Robin~Wydaeghe%
  \thanks{R. Wydaeghe is with Ghent University and imec, Ghent,
  Belgium (e-mail: robin.wydaeghe@ugent.be).}}

\markboth{IEEE Antennas and Wireless Propagation Letters,
Vol.~XX, No.~X, Month~2026}%
{Wydaeghe: Closed-form 10\,g cube peak SAR}

\maketitle

\begin{abstract}
The peak spatial-average specific absorption rate over a 10\,g cube
of tissue, prescribed by IEEE C95.3 and IEC/IEEE 62704 for
compliance below 6\,GHz and used as a sanity bound above 6\,GHz, is
typically computed by full-wave finite-difference time-domain
simulation. The computation is identified here with a closed-form
expression for any frequency where the body is opaque to the wave.
Starting from the local surface absorbed power density
$\Sab(\rr) = \Sinc T_0 \pospart{\nhat \cdot (-\khat)}$ and the
exponential depth profile of a lossy half-space, the three-dimensional
volumetric average reduces to a surface integral. For the
axis-aligned cube specified by the standard, an energy-conservation
argument gives the cube peak SAR in closed form:
$\langle\mathrm{SAR}\rangle_{\mathrm{cube}}
= \Sinc T_0 (|k_x|+|k_y|+|k_z|)/(\rhom L)$, with
$L=(m/\rhom)^{1/3} = 21.5\,$mm. The expression is exact in the
thin-skin limit ($\alpha L \gg 1$, satisfied above 6\,GHz to better
than $10^{-19}$ on skin) under the half-space and angle-independent
transmission approximations standard in millimetre-wave dosimetry.
Direct numerical integration on a tilted body half-space and on a
hemispheric sweep of incidence directions matches the closed form
to machine precision. Compliance evaluation reduces to one scalar
multiplication per surface point.
\end{abstract}

\begin{IEEEkeywords}
Specific absorption rate, peak spatial-average SAR, ICNIRP
compliance, electromagnetic dosimetry, millimetre wave,
finite-difference time-domain.
\end{IEEEkeywords}

\IEEEpeerreviewmaketitle


\section{Introduction}\label{sec:intro}
\IEEEPARstart{F}{ifth}-generation and emerging sixth-generation
wireless systems above 6\,GHz produce concentrated near-surface
energy deposition on exposed tissue, and regulators have responded
by changing the basic restriction. ICNIRP 2020 and IEC/IEEE 63195
specify the absorbed power density $\Sab$ averaged over a
$4\,\mathrm{cm}^2$ surface patch as the compliance metric above
6\,GHz, while retaining the peak spatial-average specific absorption
rate over a 10\,g volume below 6\,GHz \cite{ICNIRP2020}. The
$4\,\mathrm{cm}^2$ patch is sized to match the face of a $21.5\,$mm
cube, the shape that IEC/IEEE 62704 prescribes for the volumetric
average \cite{ICNIRP2020}. A regulatory campaign therefore traverses
the 6\,GHz boundary in both directions, and the cube quantity stays
relevant above 6\,GHz as a sanity bound, as input to combined
limits, and as the historical anchor against which all millimetre
wave dosimetry is calibrated \cite{Hirata2021,Samaras2019}.

The standardised computation is driven by full-wave simulation. The
phantom is voxelised on a Cartesian grid. An axis-aligned cube is
expanded around each tissue voxel until the enclosed mass is within
$10^{-4}$ of $10\,$g, and the maximum mass-averaged SAR over all
valid seeds is the cube peak \cite{C953,62704-1}. At millimetre-wave
frequencies the skin depth is $\delta \sim 0.3$ to $0.5\,$mm, so the
mesh must resolve a thin absorption layer over a body span of order
one metre. The cell count crosses $10^{12}$ per plane-wave
configuration, and a campaign sweeping frequency, posture, and
incidence direction takes weeks on GPU clusters.

A surface law for $\Sab$ on biological tissue at millimetre wave is
established in a sister paper \cite{WydaegheTAP_A}: at any frequency
where the body is opaque to the wave,
\begin{equation}\label{eq:surface-law}
  \Sab(\rr) = \Sinc \, T_0 \, \pospart{\nhat(\rr) \cdot (-\khat)},
\end{equation}
where $T_0$ is the polarisation-degenerate Fresnel power
transmission at normal incidence, computed in closed form from the
tissue refractive index, and $\pospart{\cdot}$ is the positive
part. The local map is a function of body geometry and incidence
direction alone. The geometric local form
\eqref{eq:surface-law} approximates the full Fresnel response to
within $5.6\,\%$ on skin at $28\,$GHz over $[0^\circ, 75^\circ]$ and
within $5\,\%$ root-mean-square across all biological tissues with
refractive index modulus above $2.5$ from $0.3$ to $100\,$GHz
\cite{WydaegheTAP_A}.

The cube peak SAR is a volumetric quantity, and its connection to
the surface law has not been written down in closed form. This
letter does so. The route is direct. The depth profile of the
internal field follows the Fresnel half-space with a single
exponential decay, the volumetric integral collapses to a surface
integral in the thin-skin limit, and energy conservation on the
axis-aligned cube returns the projected-area factor
$|k_x|+|k_y|+|k_z|$ for any incidence direction. The result is
\begin{equation}\label{eq:headline}
  \boxed{
  \langle \mathrm{SAR} \rangle_{\mathrm{cube}}
  = \frac{\Sinc \, T_0}{\rhom \, L} \,
    \bigl(|k_x| + |k_y| + |k_z|\bigr),
  }
\end{equation}
with $L = (m / \rhom)^{1/3} = 21.5\,$mm at $\rhom = 1000\,$kg/m$^3$
and $m = 10\,$g. The cube's axis-aligned geometry, the tilt of the
body surface relative to the cube, and the depth decay all collapse
into the single Cauchy projected-area factor
$|k_x|+|k_y|+|k_z|$ \cite{Cauchy1841}.

The closed form serves any compliance pipeline that already produces
a surface map of $\Sab$. The pipeline is complete with one scalar
multiplication per cube seed. Section \ref{sec:setup} fixes
notation. Section \ref{sec:reduction} performs the volume-to-surface
reduction and derives \eqref{eq:headline}. Section \ref{sec:val}
validates the formula against direct numerical integration on a
tilted body half-space across a hemisphere of incidence directions.
Section \ref{sec:disc} states the regime of validity and points to
the relaxed-geometry, sphere, and complex-anatomy extensions
treated in a separate companion paper.


\section{Setup and notation}\label{sec:setup}

A monochromatic plane wave with propagation direction $\khat$ and
incident power density $\Sinc$ illuminates a body whose surface
$\Sigma$ is locally flat on the wavelength scale. At a surface point
$\rr$, the outward unit normal is $\nhat(\rr)$ and the local
incidence cosine is
$\mu(\rr) = \nhat(\rr) \cdot (-\khat) = \cos \theta_i(\rr)$. Tissue
is modelled as a lossy half-space with complex refractive index
$\ntilde = n - j \kappa$ and field decay rate
$\alpha = k_0 \kappa$ at normal incidence. For skin at $28\,$GHz,
$\ntilde = 4.49 - 1.79 j$, $T_0 = 0.539$, and
$\alpha = 1047\,$m$^{-1}$, giving an SAR $e$-folding depth
$1/(2\alpha) = 0.48\,$mm.

The internal field intensity decays exponentially with depth $z$
along the inward normal:
\begin{equation}\label{eq:depth-decay}
  |\mathbf{E}(\rr, z)|^2
    = \frac{4\alpha}{\sigma}\, \Sab(\rr)\, e^{-2 \alpha z},
\end{equation}
which follows from integrating the ohmic loss density
$\frac{\sigma}{2}|\mathbf{E}|^2$ over depth and identifying the
result with the surface absorbed power density. The local SAR is
\begin{equation}\label{eq:sar-depth}
  \mathrm{SAR}(\rr, z)
  = \frac{\sigma}{2 \rhom}\,|\mathbf{E}(\rr, z)|^2
  = \frac{2 \alpha}{\rhom}\, \Sab(\rr)\, e^{-2 \alpha z}.
\end{equation}
The angular dependence of $\alpha$ is at most $2.2\,\%$ across
$\theta_i \in [0^\circ, 85^\circ]$ for skin at $28\,$GHz because
$|\ntilde| \gg 1$ compresses the refraction angle to within
$10^\circ$ of normal \cite{WydaegheTAP_A}. We use the
normal-incidence value throughout.

The peak spatial-average SAR over a 10\,g region is
\begin{equation}\label{eq:psSAR-def}
  \mathrm{psSAR}_{10\mathrm{g}}
  = \sup_{\Omega}\, \frac{1}{m}
    \int_{\Omega}\, \mathrm{SAR}(\rr,z)\, \rhom\, \diff V,
\end{equation}
with $m = \rhom \mathrm{Vol}(\Omega) = 10\,$g. Under
ICNIRP 2020, $\Omega$ is restricted to a cube of side
$L = (m/\rhom)^{1/3}$, axis-aligned with the global
Cartesian grid that hosts the voxel mesh. The supremum runs over
cube positions. The cube does not rotate to track the surface
\cite{C953}.


\section{Cube reduction}\label{sec:reduction}

\begin{figure*}[!t]
  \centering
  \includegraphics[width=0.92\textwidth]{psSAR_validation.pdf}
  \caption{Validation of the closed-form cube
  psSAR$_{10\mathrm{g}}$ for skin at $28\,$GHz
  ($\Sinc = 10\,$W/m$^2$, $T_0 = 0.539$, $\rhom = 1000\,$kg/m$^3$,
  $L = 21.5\,$mm). Panel (a): tilt cancellation, $\khat =
  -\hat{\mathbf{z}}$, body normal rotated by $\gamma$ in the $xz$
  plane. $\langle\Sab\rangle$ falls as $\cos\gamma$, the patch
  area inside the cube grows as $1/\cos\gamma$, and the product is
  the constant $\Sinc T_0/(\rhom L)$. Panel (b): face-flush
  configuration ($\nhat = \hat{\mathbf{z}}$), closed form contracts
  to $|k_z| = \cos\theta$ and overlaps the surface integral at
  $\phi = 0$ and $\phi = 45^\circ$ to machine precision.}
  \label{fig:val}
\end{figure*}

The cube $\mathcal{C} = [0,L]^3$ is placed near the body surface.
Tissue occupies the half-space below the surface plane
$\nhat \cdot \rr = D$. Absorbed power inside the cube is
\begin{equation}\label{eq:Pabs-volume}
  P_{\mathrm{abs}}
  = \int_{\mathcal{C}\,\cap\,\text{tissue}}\,
    \mathrm{SAR}(\rr,z)\, \rhom\, \diff V.
\end{equation}
Two routes evaluate this integral. The first integrates depth first;
the second uses energy conservation on the cube as a whole.

\subsection{Depth integral and surface collapse}
At any tissue point $(\rr, z)$ inside the cube, integrating
\eqref{eq:sar-depth} over depth from the surface plane down to the
furthest tissue point in the cube gives, in the thin-skin limit
$\alpha L \gg 1$,
\begin{equation}\label{eq:depth-collapse}
  \int_0^{d(\rr)}\, \mathrm{SAR}(\rr, z)\, \rhom\, \diff z
  = \Sab(\rr)\, \bigl(1 - e^{-2 \alpha d(\rr)}\bigr)
  \approx \Sab(\rr).
\end{equation}
At $28\,$GHz, $\alpha L \approx 45$, and the residual is below
$e^{-90} \approx 10^{-39}$. The volumetric integral
\eqref{eq:Pabs-volume} therefore reduces to a surface integral over
the body-surface patch intercepted by the cube:
\begin{equation}\label{eq:Pabs-surface}
  P_{\mathrm{abs}} \approx
  \int_{\Sigma\, \cap\, \mathcal{C}}\, \Sab(\rr)\, \diff A.
\end{equation}
The thin-skin approximation is conservative: the residual is at the
exponential-suppression scale and contributes nothing measurable.
The depth coordinate has dropped out.

\subsection{Energy conservation on the cube}
The thin-skin limit also says that all power crossing into tissue
through the body surface is absorbed within the first few skin
depths. The total absorbed power inside the cube therefore equals
the net inward Poynting flux through the air-facing portion of the
cube boundary. For an axis-aligned cube of side $L$ exposed to a
plane wave with direction $\khat$, the projected area perpendicular
to $\khat$ is the Cauchy projection of the cube,
\begin{equation}\label{eq:Aproj-cube}
  \Aproj(\khat) = L^2\, \bigl(|k_x| + |k_y| + |k_z|\bigr),
\end{equation}
which equals $L^2$ for incidence along a cube axis and
$\sqrt{3}\, L^2$ along the space diagonal. Multiplying by
$\Sinc T_0$, the absorbed power is
\begin{equation}\label{eq:Pabs-energy}
  P_{\mathrm{abs}} = \Sinc \, T_0 \, \Aproj(\khat)
  = \Sinc \, T_0 \, L^2 \,
    \bigl(|k_x|+|k_y|+|k_z|\bigr).
\end{equation}
The mass-averaged SAR is $P_{\mathrm{abs}} / m$ with
$m = \rhom L^3$, which gives \eqref{eq:headline}.

\begin{remark}[Tilt cancellation]\label{rem:tilt}
For a body surface tilted by angle $\gamma$ from the cube
$z$-axis and a wave along $\khat = -\hat{\mathbf{z}}$, the surface
patch inside the cube has area $L^2/\cos\gamma$ (tilted plane cutting
a square prism), and the local incidence cosine is $\cos\gamma$.
The surface integral \eqref{eq:Pabs-surface} returns
$\Sinc T_0 \cos\gamma \cdot L^2/\cos\gamma = \Sinc T_0 L^2$,
independent of $\gamma$. The closed form
\eqref{eq:headline} returns the same value because
$|k_x|+|k_y|+|k_z| = 1$ for axis-aligned $\khat$. The
$\cos\gamma$ in $\Sab$ exactly cancels the $1/\cos\gamma$ in
patch area.
\end{remark}

\subsection{Discrete form on a mesh}
For a body discretised as a triangulated mesh with $M$ triangles of
areas $a_j$, centroids $\rr_j$, and surface normals $\nhat_j$,
\eqref{eq:Pabs-surface} becomes
\begin{equation}\label{eq:discrete}
  \mathrm{psSAR}_{10\mathrm{g}}
  \approx \frac{1}{m}\, \max_{\mathcal{C}}\,
    \sum_{j:\rr_j \in \mathcal{C}}
    \Sinc T_0\, \pospart{\nhat_j \cdot (-\khat)}\, a_j,
\end{equation}
where the maximum runs over cube placements anchored at mesh
vertices. The closed form \eqref{eq:headline} replaces the sum by a
single scalar multiplication when the body fills the cube
cross-section perpendicular to $\khat$, which is the case for the
torso, the limbs, and the head at $L = 21.5\,$mm.


\section{Numerical validation}\label{sec:val}

The closed form \eqref{eq:headline} is tested against direct
numerical integration of \eqref{eq:Pabs-surface} on a flat body
half-space. The body surface is the plane $\nhat \cdot \rr = 0$
through the cube centre, parametrised in two cube-face coordinates
and reconstructed on a $250 \times 250$ grid. The inside-cube mask
restricts the patch to $\Sigma \cap \mathcal{C}$.
\Cref{fig:val} reports the result.

Panel (a) sweeps the body tilt $\gamma$ at axis-aligned $\khat$.
The local $\Sab$ falls as $\cos\gamma$ while the patch area
$L^2/\cos\gamma$ grows by the same factor. Their product is the
constant $\Sinc T_0 / (\rhom L)$, predicted by both the closed form
and the surface integral to within $2 \times 10^{-14}$ relative.
This is the regulatory regime: an axis-aligned cube near a tilted
body surface returns the same psSAR independent of tilt.

Panel (b) sweeps the wave direction over the upper hemisphere with
the body face-flush ($\nhat = \hat{\mathbf{z}}$). The closed form
\eqref{eq:headline} contracts to $\Sinc T_0 |k_z| / (\rhom L)$
because only the top cube face is air-facing. The surface integrals
at azimuth $\phi = 0$ and $\phi = 45^\circ$ overlap the closed form
across the entire $\theta \in [0^\circ, 89^\circ]$ range, with
maximum relative error $3 \times 10^{-14}$, set by floating-point
arithmetic. The dependence on $\phi$ vanishes by the planar symmetry
of the half-space.

For a body geometry that fills the cube cross-section perpendicular
to $\khat$, multiple cube faces face air and the full
$|k_x|+|k_y|+|k_z|$ factor is realised. This is the case for
contiguous body parts at the cube scale of $21.5\,$mm such as the
torso and the limbs, where the body occupies the entire wavefront
silhouette of the cube. For thin or folded anatomy, the surface
integral form \eqref{eq:Pabs-surface} stays exact and the closed
form is an upper bound \cite{WydaegheTAP_A}.


\section{Discussion}\label{sec:disc}

The closed form \eqref{eq:headline} is exact under the
approximations of \cite{WydaegheTAP_A}: a planar tissue interface,
the universal-decay-rate approximation on $\alpha$, and the
pseudo-Brewster compensation that makes $T_0$ the angle-independent
transmission. The thin-skin requirement $\alpha L \gg 1$ is the
weakest link below 6\,GHz, where $\alpha L$ falls to order unity for
muscle and to order $10^{-1}$ for fat \cite{Gabriel1996}. The
formula is therefore tagged as an above-6\,GHz construction.
\Cref{tab:summary} collects the numerical values for skin at
$28\,$GHz.

\begin{table}[ht]
  \centering
  \caption{Numerical values for skin at $28\,$GHz, $m = 10\,$g,
  $\rhom = 1000\,$kg/m$^3$.}
  \label{tab:summary}
  \begin{tabular}{@{}lll@{}}
    \toprule
    Quantity & Symbol & Value \\
    \midrule
    Refractive index & $\ntilde$ & $4.49 - 1.79 j$ \\
    Normal-incidence transmission & $T_0$ & $0.539$ \\
    Field decay rate & $\alpha$ & $1047\,$m$^{-1}$ \\
    Cube side length & $L$ & $21.5\,$mm \\
    Thin-skin parameter & $\alpha L$ & $45$ \\
    \addlinespace
    psSAR at $\Sinc = 20\,$W/m$^2$, axis-aligned & & $0.93\,$W/kg \\
    psSAR at $\Sinc = 20\,$W/m$^2$, diagonal & & $1.61\,$W/kg \\
    \bottomrule
  \end{tabular}
\end{table}

The space-diagonal value is the worst case in the sense of the
energy-conservation upper bound, attained when the body geometry
fills the cube cross-section perpendicular to $\khat$ and the wave
illuminates the cube along its space diagonal. For an axis-aligned
incidence on a flat body, the psSAR is the smaller value
$\Sinc T_0 / (\rhom L)$. ICNIRP's reference level for the basic
restriction at $20\,$W/m$^2$ corresponds to $0.93\,$W/kg in the
axis-aligned case, comfortably below the $2\,$W/kg basic restriction
applicable below 6\,GHz \cite{ICNIRP2020}.

The closed form fits inside a broader closed-form exposure
pipeline. The whole-body absorbed power on a body in a multipath
environment admits the direction-averaged identity of
\cite{WydaegheTAP_B}, and the coherent beamforming exposure operator
on multi-antenna millimetre-wave arrays is given in
\cite{WydaegheTWC}. With \eqref{eq:headline}, the local volumetric
quantity is computed in the same closed-form construction.
Compliance evaluation across the three regulatory metrics reduces to
surface operations on $\Sab$ plus a small number of scalar
multiplications.

Three extensions are deferred to a longer paper. First, the
relaxation of the cube prescription to an arbitrary connected
surface patch yields a one-parameter optimisation in which the
depth integral is encoded by the function $h(u) = (1-e^{-u})/u$.
The result upper-bounds the cube psSAR and recovers it when the
patch is square and matched to the cube face. Second, on a sphere
the relaxed optimisation reduces to a single transcendental
equation parametrised by a dimensionless ratio of skin depth to
body radius. Third, complex anatomy with folded tissue surfaces
within a single cube (the ear pinna and the finger are the
canonical cases) requires the surface integral form
\eqref{eq:Pabs-surface} over all illuminated patches inside the
cube. The closed form is a strict upper bound there.


\section{Conclusion}\label{sec:conc}

The peak spatial-average SAR over the IEEE/IEC 10\,g cube reduces,
above 6\,GHz, to the closed form
$\langle\mathrm{SAR}\rangle_{\mathrm{cube}}
= \Sinc T_0 (|k_x| + |k_y| + |k_z|) / (\rhom L)$. The cube's
axis-aligned geometry, the tilt of the body surface, and the depth
decay all collapse into the Cauchy projected area of the cube. Any
compliance pipeline that produces a surface map of $\Sab$ obtains
the cube quantity from one scalar multiplication per seed point.
Direct integration on a tilted body half-space and on a hemispheric
sweep of incidence directions matches the formula to machine
precision.


\bibliographystyle{IEEEtran}
\bibliography{refs}

\end{document}
